

\section{Human Evaluation}
\label{sec:abtesting}
The A/B evaluation in our study was conducted by the authors, where a human judge was presented with an input, task instruction, and two candidate outputs generated by the baseline method and \ours. The setup was blind, i.e., the judges did not know which outputs were generated by which method. The judge was then asked to select the output that is better aligned with the task instruction. 
For tasks that involve A/B evaluation, we calculate the relative improvement as the percentage increase in preference rate. The preference rate represents the proportion of times annotators selected the output produced by \ours over the output from the baseline method.
\Cref{table:ab_eval} shows the results.


\begin{table}[h]
\centering
\begin{tabular}{lccc}
\toprule
\textbf{Task} & \textbf{\ours (\%)} & \textbf{Direct (\%)} & \textbf{Either (\%)} \\
\midrule
Sentiment Transfer & 75.00 & 21.43 & 3.57 \\
Acronym Generation & 44.59 & 12.16 & 43.24 \\
Response Generation & 47.58 & 19.66 & 32.76 \\
\bottomrule
\end{tabular}
\caption{Relative improvement of \ours in A/B evaluations across different tasks. The values represent normalized preferences, which correspond to the proportion of times the output generated by \ours was selected as better aligned with the task instruction over the baseline method. The evaluation was conducted for 150 examples for each dataset. The judges were not aware of the method that generated each sample.}
\label{table:ab_eval}
\end{table}
